\documentclass[10pt,a4paper]{amsart}
\usepackage[margin=19mm]{geometry}
\usepackage{amsmath,amssymb,mathtools}
\usepackage[scr=rsfs,cal=euler]{mathalfa}
\usepackage{xfrac}
\usepackage[unicode,hidelinks,bookmarks=false]{hyperref}
\newcommand{\ie}{i.e.\ }
\newcommand{\cf}{cf.\ }
\newcommand{\resp}{respectively\ }
\newcommand{\Subsection}{\subsection}
\providecommand{\nobreakdash}{\nobreak\hbox{-}}
\setlength{\emergencystretch}{2em}
\multlinegap0pt
\usepackage{bm}
\usepackage{bbm}
\usepackage{dsfont}
\usepackage{xspace}
\usepackage{cancel}
\usepackage{mathtools}
\usepackage{amsthm}
\makeatletter
\@ifundefined{thm}{\newtheorem{thm}{Theorem}}{}
\@ifundefined{lemma}{\newtheorem{lemma}[thm]{Lemma}}{}
\makeatother
\newcommand{\N}{\mathbb{N}}
\newcommand{\RM}{\mathcal{M}[G;I,\Lambda;P]}
\newcommand{\Z}{\mathbb{Z}}
\title[Comparing Numbers of Diagonal Subsemigroups and Congruences ]{Comparing Numbers of Diagonal Subsemigroups and Congruences for Semigroups}
\author{Callum Barber, Nik Ru\v{s}kuc}
\date{Source: arXiv:2602.17286v1 (CC BY 4.0)}
\begin{document}
\maketitle

\noindent\emph{Source and licence.} Comparing Numbers of Diagonal Subsemigroups and Congruences for Semigroups. Version arXiv:2602.17286v1 (CC BY 4.0), distributed under the Creative Commons Attribution 4.0 International licence, as recorded in the arXiv licence metadata for this exact version and shown on the abstract page https://arxiv.org/abs/2602.17286v1. Full rights notice: licenses/arxiv-2602.17286-CC-BY-4.0.txt.\par\smallskip

\noindent\textit{Editorial scope.} Counted unit: the Thm whose source label is \texttt{main theorem section 3} in the article (source0.tex lines 554--562, source label \texttt{main theorem section 3}), delivered together with the article's own complete original proof of it (source0.tex lines 628--767, one \texttt{proof} environment of 140 lines). No mathematics is written, repaired or re-derived: statement and proof are the article's own text line for line. The proof is detached from the statement in the source: the article states the result at lines 554--562 and prints its proof at lines 628--767, and the proof environment's own header names the statement, so the association is the article's own. Required context retained uncounted: motivation theorem (lines 188--191); gen implies linked (lines 505--511); DSC coef form for Rees matrix (lines 534--537); N = G case (lines 542--545); elementary fraction result (lines 577--584); Bell number inequality (lines 586--593). Every definition, display and label of those lines is retained unchanged. Omitted from this package and not redistributed: 1--187, 192--504, 512--533, 538--541, 546--553, 563--576, 585--585, 594--627, 768--904. Nothing inside a retained excerpt is omitted or altered: every retained body is delivered exactly as the article's lines are written, so no omission receipt of any kind accompanies this package. Citations: the retained excerpts carry no citation command at all, so no bibliography entry of the article is transcribed and no reference is redistributed. Reference adaptation: none was needed and none was made; every reference inside the delivered unit resolves inside it, so no unresolved reference or citation is produced. Printed slips kept literal: no printed word, symbol, hypothesis or number of the counted statement or of its proof was altered. Layout and class: the article's own class is \texttt{amsart} with packages including amsmath, amssymb, hyperref, amsfonts, mathtools, theoremref, amsthm, xcolor, inputenc, array, enumitem, babel, comment, multicol; the wrapper is the lane's pinned \texttt{amsart} editorial wrapper at 10pt on A4, which restates the theorem-like environments the retained text needs and adds the article's own macro definitions that the retained text needs (\texttt{\textbackslash N}, \texttt{\textbackslash RM}, \texttt{\textbackslash Z}), with extra wrapper packages (bm, bbm, dsfont, xspace, cancel, mathtools). The delivered numbering is the unit's own. Assets: no retained span contains a figure environment, a graphics-inclusion command, a PDF-page inclusion, a table or a TikZ picture; no image file is copied into this package, the package declares no asset dependency and no image file is redistributed with it. Licence: version arXiv:2602.17286v1 of this article is distributed under the Creative Commons Attribution 4.0 International licence, as recorded in the arXiv licence metadata for this exact version and shown on the abstract page https://arxiv.org/abs/2602.17286v1. A full rights notice is delivered at licenses/arxiv-2602.17286-CC-BY-4.0.txt. Characters: every retained excerpt is pure ASCII, so the delivered engine, XeLaTeX with its default Latin Modern text font, reports no missing character.
\begin{thm}
    \label{motivation theorem}
    Let $S$ be a finite semigroup. Then $S$ is DSC  if and only if $S$ is a group.
\end{thm}

\begin{lemma}
    \label{gen implies linked}
    Let $N$ be a normal subgroup of $G$.
    Suppose we have a set of relations $\mathcal{R}_I$ on $I$ such that for each $\mathcal{S} \in \mathcal{R}_I$ the triple $(N,\mathcal{S},\Delta_\Lambda)$ is linked.
    Then $(N,\langle \mathcal{R}_I \rangle,\Delta_\Lambda) $ is a linked triple.
    Dually if $\mathcal{R}_\Lambda$ is a set of relations on $\Lambda$ such that for each $\mathcal{T} \in \mathcal{R}_\lambda$ the triple $(N,\Delta_I,\mathcal{T})$ is linked, then $(N,\Delta_I,\langle \mathcal{R}_\Lambda\rangle)$ is also linked.
\end{lemma}

\begin{equation}
    \label{DSC coef form for Rees matrix}
    \chi(\RM) = \frac{|E_{I,1}||E_{\Lambda ,1}| + \dots + |E_{I,n}||E_{\Lambda,n}|}{|R_{I,1}||R_{\Lambda ,1}| + \dots +|R_{I,n}||R_{\Lambda,n}|}.
\end{equation}

\begin{equation}
    \label{N = G case}
    |E_{I,1}| = B(|I|) ,\quad |E_{\Lambda,1}| = B(|\Lambda|),\quad |R_{I,1}| = 2^{|I|^2 - |I|}, \quad|R_{\Lambda,1}| = 2^{|\Lambda|^2 - |\Lambda|}
\end{equation}

\begin{lemma}
    \label{elementary fraction result}
    Let $a_1 , \dots , a_n , b_1 , \dots , b_n \in \N$, with $\frac{a_1}{b_1} \leq \frac{a_i}{b_i}$ for all $i$.
    Then
    \begin{align*}
        \frac{a_1}{b_1} \leq \frac{a_1 + a_2 + \dots + a_n}{b_1 + b_2 + \dots + b_n}. \tag*{\qed}
    \end{align*}
\end{lemma}

\begin{lemma}
    \label{Bell number inequality}
    Let $n , r , k_1 , \dots , k_r \in \N$ be such that $n = k_1 + \dots + k_r$.
    Then:
    \begin{equation*}
        \frac{B(n)}{2^{n^2 - n}} \leq \frac{B(k_1)}{2^{k_1^2 - k_1}} \frac{B(k_2)}{2^{k_2^2 - k_2}} \cdots \frac{B(k_r)}{2^{k_r^2 - k_r}}.
    \end{equation*}
\end{lemma}

\begin{thm}
    \label{main theorem section 3}
    Let $I$ and $\Lambda$ be sets of size $a,b>1$ respectively.
    Then for any group $G$ and $\Lambda \times I$ matrix $P$ with entries from $G$ we have that:
    \begin{equation*}
        \frac{B(a)B(b)}{2^{a^2 - a}2^{b^2-b}} \leq \chi(\RM) < 1.
    \end{equation*}
    Furthermore for any rational $\alpha \in \left[ \frac{B(a)B(b)}{2^{a^2 - a}2^{b^2-b}} , 1\right)$, there is a group $G$ and a matrix $P$ with $\chi(\RM) = \alpha$.
\end{thm}

\begin{proof}[Proof of Theorem \ref{main theorem section 3}.]
    Recall the sets:
    \begin{align*}
    E_{I,k} &= \{ \mathcal{S} : (N_k,\mathcal{S},\Delta_\Lambda) \text{ is a linked equivalence triple} \}, \\
    E_{\Lambda,k} &= \{ \mathcal{T} : (N_k,\Delta_I,\mathcal{T}) \text{ is a linked equivalence triple} \}, \\
    R_{I,k} &= \{ \mathcal{S} : (N_k,\mathcal{S},\Delta_\Lambda) \text{ is a linked reflexive triple} \}, \\
    R_{\Lambda,k} &= \{ \mathcal{T} : (N_k,\Delta_I,\mathcal{T}) \text{ is a linked reflexive triple} \}.
    \end{align*}
    Let $\sigma_k = \langle E_{I,k} \rangle$. By Lemma \ref{gen implies linked} we have that $\sigma_k \in E_{I,k}$.
    The equivalence $\sigma_k$ contains every element or $E_{I,k}$ and it is easy to see that for any equivalence $\sigma^\prime \subseteq \sigma_k$ the triple $(N_k,\sigma^\prime,\Delta_\Lambda)$ is linked. 
    And so we have the following:
    \begin{equation*}
        E_{I,k} = \{ \mathcal{S} \subseteq \sigma_k : \mathcal{S} \text{ is an equivalence relation} \}.
    \end{equation*}
    Now consider $\langle R_{I,k} \rangle$. Clearly $\sigma_k \in R_{I,k}$ and so $\sigma_k \subseteq \langle R_{I,k} \rangle$.
    Furthermore, for any $\mathcal{S} \in R_{I,k}$ we have $\mathcal{S} \subseteq \langle \mathcal{S} \rangle \subseteq \langle E_{I,k} \rangle = \sigma_k$.
    Hence $\sigma_k = \langle R_{I,k} \rangle$ and:
    \begin{equation*}
        R_{I,k} = \{\mathcal{S} \subseteq \sigma_k : \mathcal{S} \text{ is a reflexive relation} \}.
    \end{equation*}
    Let $I_1 , \dots , I_l$ be the equivalence classes of $\sigma_k$, with sizes $a_1 ,\ldots ,a_l$ respectively.
    Any equivalence relation on $I$ contained in $\sigma_k$ is just a union of equivalence relations on $I_1 , \dots ,I_l$.
    As such there are $B(a_1) \dots B(a_l)$ such relations.
    And similarly any reflexive relation on $X$ contained in $\rho$ is a union of reflexive relations on $I_1 , \dots , I_l$.
    And so there is $2^{a_1^2 - a_1} \dots 2^{a_l^2 - a_l}$ such relations.
    Therefore:
    \begin{equation}
        \label{I fraction form}
        \frac{|E_{I,k}|}{|R_{I,k}|} = \frac{B(a_1) \dots B(a_l)}{2^{a_1^2 - a_1} \dots 2^{a_l^2 - a_l}}
    \end{equation}
    where $a_1 + \dots + a_l = a$.
    Dually we have relations $\tau_k$ on $\Lambda$ such that:
    \begin{align*}
        E_{\Lambda,k} & = \{\mathcal{T} \subseteq \tau_k : \mathcal{T} \text{ is an equivalence relation} \}, \\
        R_{\Lambda,k} & = \{\mathcal{T} \subseteq \tau_k : \mathcal{T} \text{ is a reflexive relation} \}.
    \end{align*}
    Let $\Lambda_1 , \ldots , \Lambda_m$ be the equivalence classes of $\tau_k$, with sizes $b_1 , \ldots , b_m$.
    Then:
    \begin{equation}
        \label{Lambda fraction form}
        \frac{|E_{\Lambda,k}|} {|R_{\Lambda,k}|} = \frac{B(b_1) \dots B(b_m)}{2^{b_1^2 - b_1} \dots 2^{b_m^2 - b_m}}
    \end{equation}
    where $b_1 + \dots + b_m = b$.
    By \eqref{N = G case}, \eqref{I fraction form}, \eqref{Lambda fraction form} and Lemma \ref{Bell number inequality} we have:
    \begin{equation*}
        \frac{|E_{I,1}||E_{\Lambda,1}| }{|R_{I,1}||R_{\Lambda,1}|} = \frac{B(a)B(b)}{2^{a^2 - a}2^{b^2 - b}} \leq \frac{B(a_1) \dots B(a_l)}{2^{a_1^2 - a_1} \dots 2^{a_l^2 - a_l}} \frac{B(b_1) \dots B(b_m)}{2^{b_1^2 - b_1} \dots 2^{b_m^2 - b_m}} = \frac{|E_{I,k}||E_{\Lambda,k}|}{|R_{I,k}||R_{\Lambda,k}|}
    \end{equation*}
    for each $k = 1 \dots , n$.
    Finally by \eqref{DSC coef form for Rees matrix} and Lemma \ref{elementary fraction result} we get:
    \begin{equation*}
        \frac{B(a) B(b)}{2^{a^2 - a}2^{b^2-b}} = \frac{|E_{I,1}||E_{\Lambda,1}| }{|R_{I,1}||R_{\Lambda,1}|} \leq  \frac{|E_{I,1}||E_{\Lambda ,1}| + \dots + |E_{I,n}||E_{\Lambda,n}|}{|R_{I,1}||R_{\Lambda ,1}| + \dots +|R_{I,n}||R_{\Lambda,n}|} =  \chi(S).
    \end{equation*}
    As $a,b >1$ it follows that $S$ is not a group, so by Theorem \ref{motivation theorem} $\chi(S) < 1$.
    This completes the proof of the the first statement in the theorem.

    For the second part we want to show that for any rational $\alpha \in \left[\frac{B(a)B(b)}{2^{a^2 - a}2^{b^2 - b}},1\right)$, there is a group $G$ and a matrix $P$ with $\chi(\RM) = \alpha$. 
    We saw in the first part that:
    \begin{equation*}
        \frac{|E_{I,k}|}{|R_{I,k}|} = \frac{B(a_1)\dots B(a_l)}{2^{a_1^2 - a_1}\dots2^{a_l^2 - a_l}} 
    \end{equation*}
    where $a_1 , \dots , a_l$ are the sizes of the equivalence classes of $\sigma_k$.
    Suppose that for some $k$ we have $\sigma_k = I \times I$; this is certainly the case when $k = 1$, as then $N_k = G$.
    In this situation we have:
    \begin{equation*}
        \frac{|E_{I,k}|}{|R_{I,k}|} = \frac{B(a)}{2^{a^2 - a}}.
    \end{equation*}
    At the other extreme we might have $\sigma_k = \Delta_I$ for some $k$; this case may or may not happen for a particular Rees matrix semigroup.
    If it does, then:
    \begin{equation*}
        \frac{|E_{I,k}|}{|R_{I,k}|}= \frac{B(1) \dots B(1)}{2^{1^2 - 1} \dots 2^{1^2 - 1}} = 1.
    \end{equation*}
    Analogous results hold for $\frac{|E_{\Lambda,k}|}{|R_{\Lambda,k}|}$.
    
    We shall show that there exists a group $G$ and $\Lambda \times I$ matrix $P$ with entries from $G$ such that:
    \begin{itemize}
        \item $ \frac{|E_{I,k}|}{|R_{I,k}|}$ can only take values $\frac{B(a)}{2^{a^2 -a}}$ or $1$,
        \item $\frac{|E_{\Lambda,k}|}{|R_{\Lambda,k}|}$ can only take values $\frac{B(b)}{2^{b^2 -b}}$ or $1$,
        \item if there are $c$ values of $k$ for which $\frac{|E_{I,k}|}{|R_{I,k}|} = \frac{B(a)}{2^{a^2 -a}}$ then $\frac{|E_{\Lambda,k}|}{|R_{\Lambda,k}|} = \frac{B(b)}{2^{b^2 - b}}$ for the same $c$ values,
        \item if there are $d$ values of $k$ for which $\frac{|E_{I,k}|}{|R_{I,k}|} = 1$ then $\frac{|E_{\Lambda,k}|}{|R_{\Lambda,k}|} = 1$ for the same $d$ values.
    \end{itemize}
    It will then follow from \eqref{DSC coef form for Rees matrix} that we will be able to write the DSC coefficient of this Rees matrix semigroup as:
    \begin{equation*}
        \chi (\RM) = \frac{cB(a)B(b) + d}{c 2^{a^2 -a}2^{b^2 - b}+d}.
    \end{equation*}
    Furthermore we will show that we can obtain any $c \in \N$ and $d \in \N_0$.

    To this end let $p$ be a prime greater than $2^{ab+1} $ and let $G = \Z_{p^k}$, the cyclic group of order $p^k$, for some arbitrary $k \in \N_0$.
    The normal subgroups of $G$ form the chain:
    \begin{equation*}
    \textbf{1} = p^k \Z_{p^k} \leq p^{k-1} \Z_{p^k} \leq \dots \leq p \Z_{p^k} \leq \Z_{p^k} = G.
    \end{equation*}
    Let $r \in [0,k]$ be arbitrary and consider the $ab$ numbers:
    \begin{equation*}
        p^r,2p^r, 2^2 p^r, \dots,2^{ab-2} p^r , 2^{ab-1} p^r.
    \end{equation*}
    Let $P$ be any $\Lambda \times I$ matrix with all of these numbers as its entries.
    We will show that any extract $q_{\lambda \mu i j}$ with $\lambda \neq \mu$ and $i \neq j$ is a generator for $p^r \Z_{p^k}$. 
    If we consider the entries of the matrix as elements of $\Z$ this is equivalent to showing that $p^r$ divides $q_{\lambda \mu i j}$ and that $p^{r+1}$ does not divide $q_{\lambda \mu i j }$. Let:
    \begin{equation*}
        p_{\lambda i} = p^r 2^s, \quad p_{\mu i} = p^r 2^t, \quad p_{\mu j} = p^r 2^u, \quad p_{\lambda j } = p^r 2^v.
    \end{equation*}
    Then by definition $q_{\lambda \mu i j} = p^r(2^s - 2^t + 2^u - 2^v)$, and so $p^r$ divides $q_{\lambda \mu i j}$. 
    Note that:
    \begin{align*}
        | q_{\lambda \mu i j}| & = p^r| 2^s - 2^t +2^u -2^v| \\
        & \leq  p^r( 2^s + 2^t +2^u +2^v) \\
        & \leq p^r(2^{ab - 1} + 2^{ab - 1} + 2^{ab - 1} + 2^{ab - 1}) = p^r2^{ab+1} < p^{r+1}.
    \end{align*}
    As each element of $P$ is distinct we have that $s,t,u$ and $v$ are distinct, without loss of generality let $s$ be the largest of these.
    Then:
    \begin{align*}
        | q_{\lambda \mu i j}| & = p^r|2^s - 2^t + 2^u -2^v|\\
        & \geq p^r(2^s - (2^t +2^u + 2^v)) \\
        & \geq p^r(2^s - (2^{s-1} + 2^{s-2} + 2^{s-3})) \\
        & > p^r(2^s - (2^{s-1} + 2^{s-2} + 2^{s-2})) = 0.
    \end{align*}
    If follows from $0 < |q_{\lambda \mu i j}| < p^{r+1}$, that $p^{r+1}$ does not divide $q_{\lambda \mu i j}$.
    Hence $q_{\lambda \mu i j}$ is a generator for $p^r\Z_{p^k}$ for any distinct $\lambda ,\mu \in \Lambda$ and $i,j \in I$.
    
    If $N$ is a normal subgroup containing $p^r \Z_{p^k}$, then $q_{\lambda \mu i j} \in N$ for any $\lambda,\mu \in \Lambda$ and $i,j \in I$.
    So for any relations $\mathcal{S}$ on $I$ and $\mathcal{T}$ on $\Lambda$ the triple $(N,\mathcal{S},\mathcal{T})$ is linked.
    There are $r+1$ such subgroups.
    Any other normal subgroup $N$ is properly contained in $p^r \Z_{p^k}$ and so does not contain any generator for $p^r\Z_{p^k}$.
    The only extracts $q_{\lambda \mu i j}$ that are not generators for $p^r \Z_{p^k}$ have either $\lambda = \mu$ or $i =j$.
    It follows that $(N,\mathcal{S}, \mathcal{T})$ is a linked triple if and only if $\mathcal{S} = \Delta_I$ and $\mathcal{T} = \Delta_\Lambda$.
    There is $k-r$ such subgroups.
    We finally have that:
    \begin{equation}
        \label{rees matrix construction}
        \chi(S) = \frac{(r+1) B(a)B(b) + k-r}{(r+1)2^{a^2 - a} 2^{b^2 - b} + k - r}.
    \end{equation}
    For any $c \in \N$ and $d \in \N_0$ we can let $k = c+d - 1 \in \N_0$ and $r = c - 1 \leq k$, $\chi(S)$ now becomes:
    \begin{equation*}
        \frac{cB(a)B(b)+ d}{c 2^{a^2 - a} 2^{b^2 - b} + d}
    \end{equation*}
    as claimed.
    To complete the proof let $\beta ,\gamma \in \N$ with $\alpha = \frac{\beta}{\gamma}\in \left[ \frac{B(a)B(b)}{2^{a^2 - a}2^{b^2 - b}} , 1\right) $.
    Letting $c = \gamma - \beta \in \N$ and $d = \beta 2^{a^2 - a} 2^{b^2 - b} - \gamma B(a) B(b) \in \N_0$
    gives $\chi(S) = \alpha$.
\end{proof}

\end{document}
